% SOURCE official published24:2--24:3. Exact logical setting and source inference diagrams.
\section*{The Basics: Types, terms, formulas, sequents, theories}
Following Church [8], we base the syntax of terms and formulas on the simply typed $\lambda$-calculus. We reserve the primitive type $o$, and assume that there is a fixed set of primitive types (also called sorts) $\pmS$ that do not contain $o$. When $\sigma\in\pmS$, we shall also use $\sigma$ as a predicate of type $\sigma\to o$.
Our first-order, multi-sorted intuitionistic logic is based on the following propositional connectives
\[\pmt:o,\quad\wedge:o\to o\to o,\quad\pmf:o,\quad\vee:o\to o\to o,\quad\supset:o\to o\to o,\]
and the following connectives involve term structure (here, $\sigma\in\pmS$)
\[=_\sigma:\sigma\to\sigma\to o,\quad\forall_\sigma:(\sigma\to o)\to o,\quad\exists_\sigma:(\sigma\to o)\to o.\]
(We shall occasionally omit writing a typing subscript on an equality or quantifier when it can be inferred from context.) We work only in a first-order subset of Church's framework by restricting the types assigned to non-logical constants as follows.
\begin{itemize}
\item $P:\sigma_1\to\cdots\to\sigma_n\to o$ is an $n$-ary predicate symbol and $\{\sigma_1,\ldots,\sigma_n\}\subseteq\pmS$ ($n\geq0$).
\item $f:\sigma_1\to\cdots\to\sigma_n\to\sigma_0$ is an $n$-ary constructor of type $\sigma_0$ and $\{\sigma_0,\ldots,\sigma_n\}\subseteq\pmS$ ($n\geq0$).
\item $\sigma:\sigma\to o$ is a unary predicate symbol reflecting the sort $\sigma$ also as a predicate.
\end{itemize}
A collection of such declarations is called a first-order signature. We usually write $\Sigma_0$ to denote the signature of non-logical constants involved in any particular logical specification. A $\Sigma_0$-formula is any closed term of type $o$, all of whose non-logical constants come from $\Sigma_0$. Note that if $\Sigma_0$ contains declarations of only predicate symbols, then a $\Sigma_0$-formula contains no occurrences of constructors, even of primitive arity.
Sequents are two-sided and denoted by $\pmseq\Sigma\Gamma B$, where the left-hand side $\Gamma$ is a set of formulas, and the right-hand side $B$ is a single formula. The signature $\Sigma$ declares the types from $\pmS$ for all free variables occurring in $\{B\}\cup\Gamma$. Following Gentzen's terminology, $\Sigma$ is best understood as declaring the eigenvariables of the sequent---variables that function as generic or ``freshly picked'' parameters. Formally, if $\Sigma_0$ is the fixed signature of non-logical constants (constructors and predicates), then all formulas in the sequent are $(\Sigma_0\cup\Sigma)$-formulas. However, since $\Sigma_0$ is typically held constant, we usually omit it and simply refer to these as $\Sigma$-formulas, implying that their free variables are strictly contained within $\Sigma$.
The usual inference rules for $\pmLJ$ are listed in Figure 1. Both the $\mathit{instan}$ and $\mathit{cut}$ rules are eliminable, in the sense that if a sequent is provable in $\pmLJ$ then it is provable in $\pmLJ$ without these rules.
Theories are finite sets of formulas. We shall often informally say that some formula, say $B$, holds, while what we generally mean is that $B$ is provable from some theory that might be left implicit.
If $C$ is a formula whose free variables are in the list $\pmbx$, we will write this formula as $B\pmbx$ where $B$ is the (closed) abstraction $\lambda\pmbx.C$. Here, $\lambda$ plays only a notational role and does not need to be seen as part of our first-order logic.
\paragraph{Figure 1. The inference rules for the $\pmLJ$ proof system (as printed).}
\[\frac{}{\pmseq\Sigma\Gamma\pmt}\ \pmt R\]
\[
\frac{\pmseq\Sigma\Gamma B\quad\pmseq\Sigma\Gamma C}{\pmseq\Sigma\Gamma{B\wedge C}}\ \wedge R\]
\[
\frac{\pmseq\Sigma\Gamma{B_i}}{\pmseq\Sigma\Gamma{B_1\vee B_2}}\ \vee R,\ i\in\{1,2\}.\]
\[\frac{\pmseq\Sigma{B,\Gamma}C}{\pmseq\Sigma\Gamma{B\supset C}}\ \supset R\]
\[
\frac{\pmseq{y:\sigma,\Sigma}\Gamma{B[y/x]}}{\pmseq\Sigma\Gamma{\forall_\sigma x.B}}\ \forall R\]
\[\frac{\Sigma\vdash t:\sigma\quad\pmseq\Sigma\Gamma{B[t/x]}}{\pmseq\Sigma\Gamma{\exists_\sigma x.B}}\ \exists R\]
\[
\frac{\pmseq\Sigma{B_i,\Gamma}C}{\pmseq\Sigma{B_1\wedge B_2,\Gamma}C}\ \wedge L_i,\ i\in\{1,2\}.\]
\[\frac{\Sigma\vdash t:\sigma\quad\pmseq\Sigma{B[t/x],\Gamma}C}{\pmseq\Sigma{\forall_\sigma x.B,\Gamma}C}\ \forall L\]
\[\frac{\pmseq\Sigma{\Gamma_1}B\quad\pmseq\Sigma{C,\Gamma_2}C}{\pmseq\Sigma{B\supset C,\Gamma_1,\Gamma_2}C}\ \supset L\]
\[
\frac{\pmseq\Sigma{B,\Gamma}C\quad\pmseq\Sigma{C,\Gamma}C}{\pmseq\Sigma{B\vee C,\Gamma}C}\ \vee L\]
\[
\frac{}{\pmseq\Sigma{\pmf,\Gamma}C}\ \pmf L\]
\[\frac{}{\pmseq\Sigma{B,\Gamma}B}\ \mathit{init}\]
\[
\frac{\pmseq\Sigma\Gamma B\quad\pmseq\Sigma{\Gamma,B}E}{\pmseq\Sigma\Gamma E}\ \mathit{cut}\]
\[\frac{\Sigma\vdash t:\sigma\quad\pmseq{\Sigma,x:\sigma}\Gamma B}{\pmseq\Sigma{\Gamma[t/x]}{B[t/x]}}\ \mathit{instan}\]
\paragraph{Figure 2. The rules to add to $\pmLJ$ to obtain the $\pmLJ^=$ proof system.}
\[\frac{}{\pmseq\Sigma\Gamma{t=t}}\ =_R\]
\[
\frac{(t\text{ and }s\text{ not unifiable})}{\pmseq\Sigma{s=t,\Gamma}C}\ =_L\]
\[\frac{(\theta=\operatorname{mgu}(t,s))\quad\pmseq{\Sigma\theta}{\Gamma\theta}{C\theta}}{\pmseq\Sigma{s=t,\Gamma}C}\ =_L\]
An invariant of the sequents considered in this paper is that, for any sequent $\pmseq\Sigma\Gamma B$, a variable declaration $x:\sigma$ appears in $\Sigma$ if and only if the atomic formula $\sigma x$ appears in $\Gamma$. This property---the typing redundancy invariant---is preserved provided that the quantifiers $\forall_\sigma x.B$ and $\exists_\sigma x.B$ are treated as abbreviations for $\forall x_\sigma.\sigma x\supset B$ and $\exists x_\sigma.\sigma x\wedge B$, respectively.
\section*{Inference rules for equality}
The sequent calculus proof rules for the introduction of equality are presented in Figure 2. In the two versions of the $=_L$ rule, unification is done with respect to the eigenvariables. The signature $\Sigma\theta$ is the result of removing from $\Sigma$ all variables in the domain of $\theta$ and adding in all variables free in a term in the domain. When these rules are added to the $\pmLJ$ proof system, we get the proof system $\pmLJ^=$. It is important to understand that equality is a logical connective and that these inference rules fix its meaning. The addition of axioms does not affect this notion of equality. Since we shall reserve the term ``atomic formula'' for formulas that are not top-level logical connectives, a formula such as $t=s$ is not an atomic formula.
A formula is equality-free if it contains no occurrences of equality.
% SOURCE published24:4--24:5. Complete author compatibility theorem and proof retained as support.
\section*{Equivalences and compatibilities}
For every primitive type $\sigma\in\pmS$, we assume that we have a binary predicate $\pmeq_\sigma$: that is, a non-logical constant of type $\sigma\to\sigma\to o$. We abbreviate by $\pmEs$ the set $\{\pmeq_\sigma\}_{\sigma\in\pmS}$ of such binary predicates and by $\pmEa$ the set of formulas that axiomatize the properties of reflexivity, symmetry, and transitivity for all the predicates in $\pmEs$. Given this axiomatization, these binary predicates denote equivalence relations.
Let $P$ be an abstraction with type $\sigma_1\to\cdots\to\sigma_n\to o$ where $n\geq0$. We say that $P$ is compatible with respect to $\pmEs$ (following terminology in, say, [7]), if the following formula holds.
\[\text{Compatibility:}\qquad\forall\pmbx\forall\pmby.\Big(\bigwedge_{i=1}^n\pmeq\,x_i\,y_i\Big)\supset P\pmbx\supset P\pmby.\]
The equivalence predicates $\pmeq_\sigma$ themselves are trivial examples of predicates that are compatible with $\pmEs$.
Let $\pmCs$ be a set of predicates and define $\pmCa$ to be the set of formulas that axiomatize the compatibility of the predicates in $\pmCs$. The following theorem proves that any abstraction over formulas whose only non-logical constants are compatible with $\pmEs$ is also similarly compatible.
The theorems in this section consider only formulas that are equality-free: this is necessary since there might be a type $\sigma\in\pmS$ such that $\pmeq_\sigma$ is different from $=_\sigma$: this would cause a problem since we intend to lift $\pmeq_\sigma$ to be equality in a certain sense. It is, however, perfectly possible for $\pmeq_\sigma$ to be the same as $=_\sigma$. Equality plays a role starting in the next section when we define determinacy.
\pmorig{Theorem}{1} Let $B$ be an abstraction of type $\sigma_1\to\cdots\to\sigma_n\to o$ ($n\geq0$) of an equality-free formula for which all non-logical symbols are in $\pmEs\cup\pmCs$. The following compatibility-like property is provable in $\pmLJ^=$:
\[\pmEa,\pmCa\vdash\forall\pmbx\pmby.\Big(\bigwedge_i\pmeq\,x_i\,y_i\Big)\supset B\pmbx\supset B\pmby.\]
\begin{proof}
Since the displayed sequent above has no occurrences of equality, it is provable in $\pmLJ^=$ if and only if the sequent
\[\pmseq{x_1,y_1,\ldots,x_n,y_n}{\pmEa,\pmCa,\pmeq\,x_1\,y_1,\ldots,\pmeq\,x_n\,y_n,B\pmbx}{B\pmby}\]
is provable in $\pmLJ$. We proceed by induction on the structure of $B\pmbx$ to show that this sequent is provable. Besides being equality-free, $B\pmbx$ will necessarily contain no term constructors. Our proof is a simple modification of the proof that the general initial rule is admissible when only atomic initial rules are available.\par
\emph{Case. $B\pmbx$ is an atom.} There are three kinds of atomic formulas possible here.
\begin{enumerate}
\item If $B\pmbx$ is of the form $\pmeq_\sigma x_i x_j$, then the proof is completed using symmetry and transitivity for $\pmeq_\sigma$.
\item If $B\pmbx$ is of the form $\sigma x_i$ then we know that $y_i$ is of the same type as $x_i$ and, given the typing redundancy of sequents, we know that $\sigma y_i\in\Gamma$.
\item The remaining case is when this atom is of the form $(p\,u_1\cdots u_i)$ for $i\geq0$, $p\in\pmCs$, and $u_1,\ldots,u_i$ is a list of variables in $\pmbx$. This case is proved directly by using the compatibility property for $p$ that appears in $\pmCa$.
\end{enumerate}
\emph{Case. $B\pmbx$ is $B_1\pmbx\supset B_2\pmbx$.} The inductive hypotheses give proofs of two sequents
\begin{gather*}
\pmseq{x_1,y_1,\ldots,x_n,y_n}{\pmEa,\pmCa,\pmeq\,x_1\,y_1,\ldots,\pmeq\,x_n\,y_n,B_1\pmby}{B_1\pmbx},\\
\pmseq{x_1,y_1,\ldots,x_n,y_n}{\pmEa,\pmCa,\pmeq\,x_1\,y_1,\ldots,\pmeq\,x_n\,y_n,B_2\pmbx}{B_2\pmby}.
\end{gather*}
The proof of the first of these also requires the symmetry property of the equivalence relations. From these two subproofs, one can immediately build a proof of
\[\pmseq{x_1,y_1,\ldots,x_n,y_n}{\pmEa,\pmCa,\pmeq\,x_1\,y_1,\ldots,\pmeq\,x_n\,y_n,B_1\pmbx\supset B_2\pmbx}{B_1\pmbx\supset B_2\pmbx}.\]
The cases for the other binary propositional connectives ($\wedge$ and $\vee$) are completely similar. The cases for $\pmt$ and $\pmf$ are immediate.\par
\emph{Case. $B\pmbx$ is $\forall u,B_0 u\pmbx$ where $B_0$ is an $n+1$ abstraction.} The inductive hypothesis yields a proof of
\[\pmseq{x_1,y_1,\ldots,x_n,y_n,u,v}{\pmEa,\pmCa,\pmeq\,x_1\,y_1,\ldots,\pmeq\,x_n\,y_n,\pmeq\,u\,v,B_0u\pmbx}{B_0v\pmby}.\]
By substituting $u$ for $v$ into this sequent, cutting with the proof of $\pmeq\,u\,u$, and introduction rules for $\forall$ on the left and right, we have a proof of
\[\pmseq{x_1,y_1,\ldots,x_n,y_n,u,v}{\pmEa,\pmCa,\pmeq\,x_1\,y_1,\ldots,\pmeq\,x_n\,y_n,\forall u.B_0u\pmbx}{\forall u.B_0u\pmby}.\]
The case for the existential quantifier is similar.
\end{proof}
It is useful at times to replace the formulas in $\pmEa$ and $\pmCa$ with synthetic rules. By employing the techniques found in [26,33], these formulas can also be turned into the following inference rules, encoding, respectively, reflexivity, symmetry, transitivity, and compatibility.
\[\frac{\pmseq{u:\sigma,\Sigma}{\pmeq_\sigma u u,\Gamma}B}{\pmseq{u:\sigma,\Sigma}\Gamma B}\]
\[
\frac{\pmseq\Sigma{\pmeq_\sigma t s,\pmeq_\sigma s t,\Gamma}B}{\pmseq\Sigma{\pmeq_\sigma t s,\Gamma}B}\]
\[\frac{\pmseq\Sigma{\pmeq_\sigma t s,\pmeq_\sigma s r,\pmeq_\sigma t r,\Gamma}B}{\pmseq\Sigma{\pmeq_\sigma t s,\pmeq_\sigma s r,\Gamma}B}\]
\[\frac{\pmseq\Sigma{\pmeq_{\sigma_1}t_1s_1,\ldots,\pmeq_{\sigma_1}t_1s_1,P\,t_1\cdots t_n,P\,s_1\cdots s_n,\Gamma}B}{\pmseq\Sigma{\pmeq_{\sigma_1}t_1s_1,\ldots,\pmeq_{\sigma_1}t_1s_1,P\,t_1\cdots t_n,\Gamma}B}\]
As shown in [26], the addition of such synthetic rules to $\pmLJ$ results in a proof system that satisfies the cut elimination theorem.
% SOURCE published24:5--24:7. Complete counted author proof-system equivalence.
\section*{Inference rules for equivalences}
If we are working only with equality-free formulas over $\pmEs\cup\pmCs$, then we can move to a new proof system in which the (possibly numerous) axioms describing compatibility and equivalence are replaced by inference rules that make the predicates in $\pmEs$ into logical connectives by giving them their own left and right introduction rules. In particular, those new introduction rules are essentially the rules for equality.
\[\frac{}{\pmseq\Sigma\Gamma{\pmeq_\sigma t t}}\ \pmeq_\sigma R\qquad
\frac{\pmseq{u,\Sigma}{\Gamma[u/v]}{B[u/v]}}{\pmseq{u,v,\Sigma}{\pmeq_\sigma u v,\Gamma}B}\ \pmeq_\sigma L\]
Given that we will be working with formulas without occurrences of constructors, the rules for introducing equivalences on the left and right are simpler than for general equality. In particular, the left introduction rule lacks the ``failure-to-unify'' case. Let $\pmLJ_{\pmEs}$ be the proof system that results from adding the above introduction rules (for every $\sigma\in\pmS$) to $\pmLJ$, and let $\pmLJ^=_{\pmEs}$ be the proof system that results from adding them to $\pmLJ^=$.
% COUNTED STATEMENT 2
\pmorig{Theorem}{2} Let $B$ be an $\pmEs\cup\pmCs$-formula without equalities. If $\pmseq\cdot{\pmEa,\pmCa}B$ has an $\pmLJ^=$ proof then $\pmseq\cdot\cdot B$ has an $\pmLJ^=_{\pmEs}$ proof.
% COUNTED PROOF 2
\begin{proof}
Assume $\pmseq\cdot{\pmEa,\pmCa}B$ has an $\pmLJ^=$ proof. Since $B$ contains no equalities, then $\pmseq\cdot{\pmEa,\pmCa}B$ has an $\pmLJ$ proof. For the purpose of this proof, we extend the $\pmLJ$ proof system to the $\pmLJ^+$ proof system by adding the synthetic rules for all the formulas in $\pmEa\cup\pmCa$: as a result, the assumptions in $\pmEa\cup\pmCa$ can be removed from the sequent. In that extended proof system, every sequent is of the form
\begin{equation}\pmseq{u_1:\sigma_1,\ldots,u_n:\sigma_n}{S,R,\Gamma}C,\tag{1}\end{equation}
where $S$ is $\{\sigma_1u_1,\ldots,\sigma_nu_n\}$, and $R$ is the set of atomic formulas of the form $\pmeq_\sigma u_i u_j$ (and we assume that $\Gamma$ does not contain such atomic formulas). For $\sigma\in\pmS$, let $\equiv_\sigma$ be the least equivalence relation on the variables in $\{u_1,\ldots,u_n\}$ that contains the relation $\{\langle u,v\rangle\mid\pmeq_\sigma u v\in R\}$. We fix some linear order on variables (say, alphabetic) and define the function $\widehat{\equiv}_\sigma$ from variables to variables (both of type $\sigma$) such that $\widehat{\equiv}_\sigma(u)=v$ if $v$ is the representative of the $\equiv_\sigma$-equivalence class containing $u$. Here, the representative is the member of the equivalence class that is least in the order relation. The function $\widehat{\equiv}(C)$ is defined to map the formula $C$ to the formula that results in substituting free variables of $C$ (of type, say, $\sigma$) with their image under $\widehat{\equiv}_\sigma$. Next, we apply $\widehat{\equiv}$ to the formulas in the sequent (1) to get the new sequent
\begin{equation}\pmseq{v_1,\ldots,v_m}{\widehat{\equiv}_\sigma(S),\widehat{\equiv}_\sigma(R),\widehat{\equiv}_\sigma(\Gamma)}{\widehat{\equiv}_\sigma(C)}.\tag{2}\end{equation}
Here, the $v_i$'s are representatives of equivalence classes and the atoms in $\widehat{\equiv}_\sigma(R)$ are all of the form $\pmeq_\sigma v_i v_i$. We shall say that the sequent (2) is associated with the sequent (1).
A simple induction on the structure of $\pmLJ^+$ proofs shows that if (1) has an $\pmLJ^+$ proof, then its associated sequent has an $\pmLJ_{\pmEs}$ proof. As a result of this induction, we know that since the sequent $\pmseq\cdot\cdot B$ has an $\pmLJ^=$ proof, then the associated sequent $\pmseq\cdot\cdot B$ has an $\pmLJ^=_{\pmEs}$ proof.
\end{proof}
% COUNTED STATEMENT 3 SAME EQUIVALENCE REVERSE DIRECTION
\pmorig{Theorem}{3} Let $B$ be an $\pmEs\cup\pmCs$-formula without equalities. If $\pmseq\cdot\cdot B$ has an $\pmLJ^=_{\pmEs}$ proof then $\pmseq\cdot{\pmEa,\pmCa}B$ has an $\pmLJ^=$ proof.
% COUNTED PROOF 3
\begin{proof}
We show how to transform a $\pmLJ^=_{\pmEs}$ proof of $\pmseq\Sigma\Gamma B$ into an $\pmLJ^=$ proof $\pmseq\Sigma{\pmEa,\pmCa,\Gamma}B$. All rules of $\pmLJ^=_{\pmEs}$ are rules of $\pmLJ^=$ except for the $\pmeq_\sigma L$ and $\pmeq_\sigma R$ rules. The $\pmeq_\sigma R$ rule is easily accounted for using the reflective rules in $\pmEa$. To account for the $\pmeq_\sigma L$ rule, we first prove that the sequents
\begin{gather*}\pmseq{u,v,\Sigma}{\pmEa,\pmCa,\pmeq_\sigma u v,B}{B[u/v]}\tag{3}\\
\pmseq{u,v,\Sigma}{\pmEa,\pmCa,\pmeq_\sigma u v,B[u/v]}B\tag{4}
\end{gather*}
have $\pmLJ^=$ proofs. We do this by induction on the structure of $B$. The case where $B$ is an atomic formula follows immediately from the axioms in $\pmEa,\pmCa$. All the other cases follow directly from the inductive hypothesis.
Now, assume that the last inference rule is the $\pmeq_\sigma L$ rule. We apply the inductive assumption to conclude that the sequent
\[\pmseq{u,\Sigma}{\pmEa,\pmCa,\Gamma[u/v]}{B[u/v]}\]
has an $\pmLJ^=$ proof. We can now apply the cut rule with (3) for every member of $\Gamma[u/v]$ and with (4) with $B[u/v]$. Cut elimination for $\pmLJ^=$ proofs then yields the desired proof of the sequent $\pmseq{u,v,\Sigma}{\pmEa,\pmCa,\pmeq_\sigma u v,\Gamma}B$.
\end{proof}
